\section{Finite Horn Closure and Incremental Support}
\label{sec:horn}

Horn support supplies the monotone base of the architecture. The immediate-consequence operator in Equation~\eqref{eq:horn-operator} is monotone because every premise inclusion that holds in (S) also holds in a superset (S'):
\begin{equation}
S\subseteq S'\Longrightarrow T_H(S)\subseteq T_H(S').
\label{eq:horn-monotone}
\end{equation}
This property is \formalized{}. It yields an ascending finite iteration rather than relying on an unbounded operational search. Since each $I_n\subseteq U$, either the chain adds a new element or it stabilizes. There can be at most $|U|$ strict additions, hence
\begin{equation}
I_{|U|}=I_{|U|+1}=T_H(I_{|U|}).
\label{eq:horn-stabilization}
\end{equation}

The fixed point is least. If $T_H(S)=S$, induction on $n$ gives $I_n\subseteq S$; therefore
\begin{equation}
C_H=I_{|U|}\subseteq S.
\label{eq:horn-least}
\end{equation}
The proof is machine checked in the finite setting. It justifies reading (C_H) as the supported fact carrier generated by the stated base and rules. It does not justify the base facts or rule encoding themselves.

A candidate (c) is well formed only if it is request-bound and supported by this closure:
\begin{equation}
\operatorname{CandidateWF}(H,c)\iff
c.r=H.r\land c.support\subseteq C_H.
\label{eq:candidate-wf}
\end{equation}
The formula is \formalized{}. It blocks a candidate from smuggling in a premise absent from the declared Horn system. It does not require that every fact in (C_H) be relevant to a candidate.

Incremental updates are restricted to additions. For $\Delta=(\Delta_F,\Delta_R)$,
\begin{equation}
H+\Delta=(F_0\cup\Delta_F,R\cup\Delta_R,U).
\label{eq:add-only-extension}
\end{equation}
The universe is fixed and the old facts and rules remain. Consequently,
\begin{equation}
T_H(S)\subseteq T_{H+\Delta}(S),
\qquad C_H\subseteq C_{H+\Delta}.
\label{eq:add-only-monotonicity}
\end{equation}
Both inclusions are \formalized{}. Retraction is excluded; a non-monotone change would require a different update theory and invalidation mechanism.

The repository specifies incremental implementation correctness extensionally:
\begin{equation}
\operatorname{IncrementalCorrect}(impl)\iff
\forall\Delta,\quad impl(\Delta)=
\operatorname{FullRecompute}(H+\Delta).
\label{eq:incremental-correct}
\end{equation}
This definition is \formalized{}, but it does not exhibit or verify a production worklist algorithm. A concrete implementation must provide a refinement witness or executable evidence. Confusing the contract with its implementation would reverse the direction of verification.

The separation between monotone support and later non-monotone acceptance is essential. New base facts or rules cannot erase previously supported facts under Equation~\eqref{eq:add-only-monotonicity}. Yet new attacks can change which arguments are accepted under an argumentation profile. The architecture therefore preserves the monotone carrier (C_H) while allowing the acceptance layer to expose disagreement rather than rewriting derivability.

This separation yields a useful diagnostic taxonomy. If an expected proposition is absent from $C_H$, the defect lies in the base facts, rule set, universe, or closure implementation. If the proposition is supported but no associated argument is accepted, the defect or disagreement lies in argument construction, attack generation, defeat policy, or extension selection. Debugging can therefore target a layer instead of treating every unfavorable result as a rule-engine failure.

The add-only restriction also defines a safe cache boundary. A cached old closure remains a subset of the new closure, but it is not necessarily equal to it. A runtime may reuse old support as a seed only if it recomputes the consequences of added facts and rules and produces the same carrier as full recomputation. Deletions, authority withdrawal, or temporal expiry do not satisfy this premise and must trigger invalidation rather than incremental addition.

From a legal perspective, monotonic support should not be confused with precedential entrenchment. A fact can remain derivable in the rule carrier while a later exception, attack, or procedural bar defeats the argument built from it. The two-layer design records both facts: the support did not disappear, and its decisional role changed.

\subsection{Proof anatomy and operational consequences}

The finite closure result depends on three distinct ingredients. Monotonicity orders the iterations, the finite universe bounds the number of strict additions, and the fixed-point argument identifies the stabilized carrier. Dropping any ingredient changes the claim. A monotone operator on an infinite universe need not stabilize after a known finite number of steps. A finite but non-monotone operator may oscillate. A stable implementation output need not be least unless it is related to the iteration from the empty carrier. The formal theorem packages these facts for the specified Horn operator; it is not a generic termination certificate for arbitrary rule engines.

This decomposition improves failure reports. If a runtime exceeds the expected bound, reviewers should ask whether the executable universe matches the formal finite carrier, whether rules introduce undeclared symbols, or whether the implementation is repeating work without changing the set. If the output stabilizes but disagrees with reference closure, the relevant issue is refinement rather than mathematical existence. If an add-only update appears to delete support, either the update violated the premise or the implementation did not realize the specified union. Each diagnosis points to a concrete witness rather than the generic statement that ``reasoning failed.''

The least-closure theorem is also deliberately neutral about legal relevance. It derives every fact licensed by the supplied Horn system, including facts that may never enter an accepted argument. A production implementation can optimize demand-driven evaluation, but the optimization must preserve the observable closure required at its interface. Performance evidence may justify the optimization; only equality with the reference semantics justifies substituting it in the assurance argument.
